\subsection{Families of Sets}
%TODO CLEANUP
Let $X$ be a set. The \textbf{power set} $\mathcal{P}(X)$, consisting of all subsets of $X$, is naturally ordered by inclusion: for $A, B \in \mathcal{P}(X)$, we define
\[
A \le B \quad \Longleftrightarrow \quad A \subseteq B.
\]

Any subset of $\mathcal{P}(X)$ inherits this order. Such a collection is called a \emph{family of sets}. Families may be specified purely set-theoretically—for example, the collection of all finite subsets of an infinite set $X$.

More commonly, families of sets arise when $X$ carries additional structure. For instance, if $X$ has an algebraic structure, such as a group, vector space, or ring, then the following collections form ordered sets under inclusion:
\begin{itemize}
    \item the set of all subgroups of a group $G$, denoted $\operatorname{Sub}(G)$, and the set of all normal subgroups, denoted $\operatorname{NSub}(G)$;
    \item the set of all subspaces of a vector space $V$, denoted $\operatorname{Sub}(V)$;
    \item the set of all subrings of a ring $R$, and the set of all ideals of $R$.
\end{itemize}

Families of sets also arise in other contexts. For example, if $(X, \mathcal{T})$ is a topological space, then the collections of open sets, closed sets, and \emph{clopen} sets (those that are both open and closed) are each ordered by inclusion.

A particularly important example is the family $\mathcal{O}(P)$ of \emph{down-sets} of a partially ordered set $P$, which will be introduced later.

\medskip

The ordered set $(\mathcal{P}(X), \subseteq)$ admits an equivalent formulation in terms of predicates. A \textbf{predicate} on $X$ is a function
\[
p : X \to \{T, F\},
\]
where $T$ and $F$ denote truth values. For example, the function $p : \mathbb{R} \to \{T, F\}$ defined by
\[
p(x) =
\begin{cases}
T & \text{if } x > 0, \\
F & \text{otherwise}
\end{cases}
\]
is a predicate on $\mathbb{R}$.

Let $\mathcal{P}^*(X)$ denote the set of all predicates on $X$. We order $\mathcal{P}^*(X)$ by implication:
\[
p \le q \quad \Longleftrightarrow \quad \{x \in X : p(x) = T\} \subseteq \{x \in X : q(x) = T\}.
\]

Define a map
\[
\Phi : \mathcal{P}^*(X) \to \mathcal{P}(X), \quad \Phi(p) := \{x \in X : p(x) = T\}.
\]
Then $\Phi$ is an order-isomorphism between $(\mathcal{P}^*(X), \le)$ and $(\mathcal{P}(X), \subseteq)$.

\begin{remark}
The identification of subsets with predicates is fundamental in logic and computer science, where sets are often represented via their characteristic functions.
\end{remark}